% ============================================================
% ch11.tex —— 第 11 章 多项式大厦：另一种数
% ============================================================
\chapter{多项式大厦：另一种数}

$x^2 + 3x + 2$ 这种东西，你初中三年从没把它当"数"看过——它是一个式子，用来代入、求值、画图。这一章请换个眼光：把多项式当成\textbf{另一种数}来看。你会发现它和整数之间的相似度高到令人起疑——而这份相似不是比喻，是\textbf{结构相同}，它能一路兑现到本书的最后一章。

\section{给多项式发户口}

先把户口办了。全体多项式（系数取自某个域，本书用分数或实数；变量 $x$；次数任意）组成的集合，配上加法（合并同类项）与乘法（逐项展开），是一个\textbf{环}——第 10 章的定义逐条可验（加法是群、乘法封闭结合有单位元"$1$"、分配律就是展开括号的规则）。记号 $F[x]$（$F$ 是系数域）。

两套运算的日常你已经熟练。现在补上缺的第三样：\textbf{除法}。初中老师对"整式除法"一笔带过，这里郑重其事算一次，用竖式：

\[
\begin{array}{r}
x^2 + x + 4 \quad\leftarrow\text{商} \\
x-1\ \big)\ \begin{array}[t]{r@{}r@{}r@{}r}
x^3 \phantom{{}+0x^2} & {}+3x & {}+2 & \\
\underline{-\,(x^3 - x^2)}\phantom{+} \\
\phantom{-(}x^2 + 3x\phantom{{}+2}\\
\underline{-\,(x^2 - x)}\phantom{+2}\\
\phantom{-(}4x + 2\\
\underline{-\,(4x - 4)}\\
\phantom{-(}6 \quad\leftarrow\text{余数}
\end{array}
\end{array}
\]

读出结果：
\[
x^3 + 3x + 2 = (x-1)\left(x^2 + x + 4\right) + 6 .
\]
对照整数的带余除法 $387 = 1\times210 + 177$：结构一模一样——"被除式 $=$ 除式 $\times$ 商 $+$ 余式"，且\textbf{余式的次数低于除式的次数}（正如整数余数小于除数）。多项式世界有自己完整的带余除法。

做除法的口诀也和整数竖式同款："拿最高次对最高次，消掉头，降一次，再来。"$x^3 \div x = x^2$，乘回去 $(x-1)x^2 = x^3 - x^2$，相减剩 $x^2 + 3x$；$x^2 \div x = x$，乘回去 $x^2 - x$，相减剩 $4x + 2$；$4x \div x = 4$，乘回去 $4x - 4$，相减剩 $6$——次数 $0$ 低于除式次数 $1$，收工。

\section{余数定理：代回就能偷看余数}

上面那道题，硬除要五分钟。现在表演一个五秒的方法：把 $x = 1$ 代进原式，
\[
f(1) = 1 + 3 + 2 = 6 .
\]
恰好是余数！巧合吗？再来一发：$x^{100}$ 除以 $(x+1)$，余数是多少？硬除是不可能的（商是一百次的多项式）。设
\[
x^{100} = (x+1)\times(\text{某个商}) + r ,
\]
其中 $r$ 是常数（余式次数必须低于 $1$，所以是常数）。把 $x = -1$ 代进去：右边带 $(x+1)$ 的项整项归零，剩下
\[
(-1)^{100} = 0 + r \quad\Longrightarrow\quad r = 1 .
\]
余数 $1$，五秒。这不是魔术，是定理：

\begin{dingli}[余数定理]
多项式 $f(x)$ 除以 $(x - a)$，余数等于 $f(a)$（把 $a$ 代回去的函数值）。
\end{dingli}
\begin{proof}
带余除法说 $f(x) = (x-a)\,q(x) + r$，$r$ 是常数。这是\textbf{恒等式}——对\textbf{一切} $x$ 成立。代入 $x = a$：左边 $f(a)$；右边 $(a-a)q(a) + r = 0\times q(a) + r = r$。所以 $f(a) = r$。证毕。
\end{proof}

两行证明，力量全在"恒等式对一切 $x$ 成立，所以挑一个好用的 $x$ 代入"——这是代数里最锋利的招式之一：\textbf{普适等式，专点化简}。（你其实早就用过这招：解方程验根时"代回去"，逻辑同源。）

立即推论（\textbf{因式定理}）：$f(a) = 0 \iff (x-a)$ 整除 $f(x)$（余数为零，即除尽）。
\textbf{"方程的根"与"因式分解"从此合流}：$x^2 - 3x + 2$ 在 $x=1$ 处取值 $0$，所以 $(x-1)$ 是因式，于是 $x^2-3x+2 = (x-1)(x-2)$——你初中用的"十字相乘法"，原理就是它；"先试根再分解"的试除法，原理还是它。更进一步：\textbf{三次方程找不到因式分解时，先试整数根}（$\pm1, \pm2, \dots$ 逐个代），找到一个 $a$ 就撕下一个 $(x-a)$，三次降二次，二次套公式。整条解方程流水线的地基，就是这两条定理。

\section{两张户口本的对照（本章主展品）}

把整数与多项式并排放好，逐项对照——这张表是全章的灵魂，也是通往篇六的大门：

\begin{center}
\begin{tabular}{lll}
\toprule
项目 & 整数世界 $\mathbb{Z}$ & 多项式世界 $F[x]$ \\
\midrule
基本对象 & $\pm1, \pm2, \pm3, \dots$ & $1,\ x,\ x^2,\ \dots$ 的组合 \\
带余除法 & $387 = 1\times210 + 177$ & $x^3{+}3x{+}2 = (x{-}1)(x^2{+}x{+}4) + 6$ \\
余数规则 & $0 \le r < $ 除数 & $\deg r < \deg$ 除式 \\
"大小" & 绝对值 & 次数（每次除法至少降一次） \\
原子 & 素数（不能拆） & 不可约多项式（不能拆） \\
唯一分解 & $360 = 2^3\cdot3^2\cdot5$ 仅此一种 & $x^4 - 1 = (x{-}1)(x{+}1)(x^2{+}1)$ 仅此一种（实系数） \\
最大公因数 & 辗转相除 & 同一套辗转相除 \\
裴蜀等式 & $ax + by = \gcd(a,b)$ & $u(x)f(x) + v(x)g(x) = d(x)$ \\
余数窥视 & —— & 余数定理：$f(a)$ \\
\bottomrule
\end{tabular}
\end{center}

"不可约"（拆不动）随\textbf{系数域}变化：$x^2 + 1$ 在实数上不可约（判别式 $-4 < 0$，没有实根，撕不下一次因式）；在复数上 $= (x-\mathrm{i})(x+\mathrm{i})$，一撕就开。同理，$x^2 - 2$ 在有理数（分数）上不可约，在实数上 $= (x-\sqrt2)(x+\sqrt2)$。\textbf{原子清单取决于你允许哪种数当系数}——第 11 章这颗种子，第 30 章数交点时发芽（交点个数也随"允不允许复数"而变），全书前后呼应的一根暗线。

多项式版的"唯一分解定理"同样成立（任何非常数多项式拆成不可约因式之积，拆法唯一不计次序与常数倍）——证明的骨架与第 3 章完全平行：先证"多项式版的欧几里得引理"（不可约式整除乘积则整除某因子，用裴蜀等式走同一条路），再逐对消去。两张户口本不仅像，\textbf{连定理的证明都逐字对应}——数学家把这种现象升级成方法论：\textbf{只要一个定理的证明只用了带余除法和四则运算，它就在两个世界同时免费成立}。这个观察的正式名字叫"欧几里得整环理论"，它是大学抽象代数的开胃菜。

\section{为什么这份相似值钱：三个直接回报}

\textbf{回报一：分式函数有了户口}。第 10 章的分数流水线开进多项式环：$\dfrac{x+1}{x^2+3}$ 这种"多项式分数"组成一个域（有理函数域），微积分第 17、18 章的积分表里大量居民来自这里。

\textbf{回报二：解方程 $=$ 分解 $=$ 数根}。$n$ 次多项式至多 $n$ 个根——用因式定理立刻可证：每个根撕下一个一次因式，$n$ 次多项式只有 $n$ 个一次因式可撕（次数守恒）。\textbf{这是第 30 章"交点乘法定律"（贝祖定理）的直系祖先}：联立两条曲线消元后得到一个 $n$ 次方程，根的个数上限就卡死了交点个数。

\textbf{回报三：拉格朗日插值——五点定一条二次曲线}。问：是否存在一条二次多项式曲线，恰好穿过指定的三个点 $(x_1,y_1), (x_2,y_2), (x_3,y_3)$（$x$ 互不相同）？答案是肯定的，且能当场造出来。诀窍是"各管一摊"（第 5 章中国剩余定理的直系子孙）：造三个"专职多项式"
\[
L_1(x) = \frac{(x-x_2)(x-x_3)}{(x_1-x_2)(x_1-x_3)}, \quad L_2, L_3\ \text{同理}，
\]
每个 $L_i$ 的性质：在自家 $x_i$ 处取值 $1$，在别人家取值 $0$（分子归零）。叠加：
\[
f(x) = y_1 L_1(x) + y_2 L_2(x) + y_3 L_3(x)
\]
三点全过 \checkmark。这个构造以拉格朗日命名，今天的飞机导航、气象插值、手机屏幕的字体重建，底层全是它。\textbf{第 5 章的魔法数、本章的专职多项式，同一招隔空击掌}——这就是"结构相同"的又一场演出。

\begin{dongshou}
\begin{itemize}
  \yisi 用五秒法求余数：$x^{99} \div (x+1)$（代入 $x=-1$：$(-1)^{99} = -1$）；$x^5 - 3x + 1 \div (x-2)$（代入：$32 - 6 + 1 = 27$）。
  \yier 完全分解 $x^3 - 6x^2 + 11x - 6$。（试 $x = 1$：$1 - 6 + 11 - 6 = 0$ \checkmark，撕下 $(x-1)$，竖式除得 $x^2 - 5x + 6 = (x-2)(x-3)$。三个根 $1, 2, 3$ 全是整数——出题人的善意。）
  \yier 对多项式执行辗转相除：求 $\gcd(x^3 - 1,\ x^2 - 1)$。（$x^3 - 1 = (x)(x^2-1) + (x - 1)$；$x^2 - 1 = (x+1)(x-1) + 0$。答案 $x - 1$——两个"三次单位根集合与二次单位根集合"的公共部分，正是 $x=1$ 这一个根。）
  \yier 用因式定理证明：$f(x)$ 是整系数多项式且 $f(a) = f(b) = 0$（$a \ne b$），则 $(x-a)(x-b)$ 整除 $f(x)$。再用它秒判：一个 $n$ 次多项式至多有 $n$ 个根（反证：若有 $n+1$ 个根，撕下 $n+1$ 个一次因式，次数超编——与"次数恰为 $n$"矛盾，除非它是零多项式）。
\end{itemize}
\end{dongshou}

\begin{shaonao}
\textbf{多项式版中国剩余定理（拉格朗日插值的孪生兄弟）。}求一个次数不超过 $2$ 的多项式 $r(x)$，使它除以 $(x-1)$ 余 $3$、除以 $(x-2)$ 余 $5$。两条思路：

思路一（待定系数）：$r(x) = (x-1)(x-2)q(x) + ax + b$（余式至多一次）。代入 $r(1) = 3$：$a + b = 3$；$r(2) = 5$：$2a + b = 5$。解得 $a = 2, b = 1$，即 $r(x) = 2x + 1$。验证：$r(1) = 3$ \checkmark，$r(2) = 5$ \checkmark。

思路二（各管一摊）：照搬第 5 章的魔法数构造——$\ell_1(x) = \dfrac{x-2}{1-2} = 2 - x$（在 $1$ 处为 $1$，在 $2$ 处为 $0$），$\ell_2(x) = x - 1$（反之），叠加 $3\ell_1 + 5\ell_2 = 3(2-x) + 5(x-1) = 2x + 1$。两种思路同一答案——\textbf{整数的同余世界与多项式的取值世界，共享同一套"分而治之"骨架}。把这条暗线记在心里：第 31 章（射影平面）与第 32 章（椭圆曲线）还会第三次、第四次踩到同一块跳板。
\end{shaonao}

\begin{chengshi}
多项式环上"唯一分解"的完整证明与整数版平行展开即可（先欧几里得引理、后消去法），我们按"结构相同"原则一笔带过——严格的版本见任何一本抽象代数教材的"欧几里得整环"一节。拉格朗日插值对 $k$ 个点推广为 $k-1$ 次多项式，公式照抄；"$x_i$ 互不相同"这个条件不可省（两点同 $x$ 不同 $y$ 时函数无定义——这正是"函数"与"曲线"的差别，篇六还会回来纠缠这个细节）。
\end{chengshi}

篇二还剩最后一个选读章：一段少年天才用"数的次数"给"尺规作图"判死刑的传奇。赶时间的读者可以直接跳去篇三，那里有无限、瞬间、以及整个宇宙最会赚钱的一个数。
